\documentclass[11pt]{article}
\usepackage[margin=1in]{geometry}
\usepackage{amsmath,amssymb,amsthm,microtype}
\usepackage{titling}
\setlength{\droptitle}{-30pt}
\usepackage[hidelinks]{hyperref}
\newtheorem{theorem}{Theorem}
\newtheorem{corollary}[theorem]{Corollary}
\newcommand{\rough}{\overline M}
\newcommand{\C}{\mathbb C}
\title{Classical arrangement counterexamples to\protect\\ summed rough-\(M\) semicontinuity}
\author{Alec Kriebel\thanks{Independent researcher. ORCID: \href{https://orcid.org/0009-0001-9320-500X}{0009-0001-9320-500X}.}}
\date{October 14, 2026}
\begin{document}
\maketitle
\begin{abstract}
The classical dual Hesse arrangement and the established deformation of an
affine line arrangement to its pencil give a delta-constant multibranch family
whose central rough invariant is \(7\), whereas the sum on every nearby fiber
is \(12\). We record the elementary corollary with explicit coefficients,
finite simultaneous normalization, and a common transverse ball. It gives a
negative answer to the unrestricted summed reading of the rough-invariant
question in Borodzik's 2009 Oberwolfach contribution. It does not address the
later explicitly unibranched parametric conjecture. We credit the earlier
construction and method and make no firstness or present-openness claim.
\end{abstract}

\paragraph{Question and provenance.}
For a reduced plane-curve singularity, let
\[
 \rough(C,p)=K\cdot(K+D),\qquad D=C'+E_{\rm red},
\]
on its minimal embedded resolution; \(K\) is the local canonical divisor
supported on the exceptional divisor. For a fiber with several singularities
we take the sum of these local values in a fixed representative.
Borodzik's report \cite[pp.~2446--2447]{OWR} defines this rough invariant and
asks about upper semicontinuity, including delta-constant deformations.
The question itself does not write a sum or impose unibranchedness; its
following fixed-ball discussion compares sums over nearby singularities.
Our conclusion concerns precisely this unrestricted summed reading, without
asserting an unstated intended scope.

The geometry is prior mathematics. De Jong--van Straten
\cite[Definition~4.2 and Example~6.6]{DvS} use the exact nine-line,
twelve-triple configuration in the delta-constant picture-deformation
framework. The scaling of affine line constants to zero is explicit in
Plamenevskaya--Starkston \cite[Example~8.7]{PS}. Borodzik
\cite[p.~578]{B} already uses contraction to an ordinary multiple point
and its rough value \(r-2\). Thus the theorem below is an explicit classical
corollary and formulation clarification; it asserts neither a new
arrangement nor a new deformation method. The inspected sources did not
establish earlier publication of this exact rough-invariant comparison;
that bounded observation does not establish priority.

\begin{theorem}\label{thm}
There is a reduced flat holomorphic family in the open unit ball in
\(\C^2\), over \(|t|<1/4\), with finite simultaneous normalization and smooth
disjoint transverse boundary circles, such that the central fiber has an
ordinary ninefold point and every nonzero fiber has exactly twelve ordinary
triple points. All singularities lie inside the ball of radius \(1/2\), and
\[
 \delta(C_0,0)=\sum_p\delta(C_t,p)=36,\qquad
 \rough(C_0,0)=7<12=\sum_p\rough(C_t,p)\quad(t\ne0).
\]
Every normalized fiber consists of nine disks. The construction restricts,
after scaling, to arbitrarily small ball representatives.
\end{theorem}

\begin{proof}
Let \(\omega\) be a primitive cube root of unity. In \(\mathbb P^2\) take
\[
 A_i:X-\omega^iY=0,\quad B_j:Y-\omega^jZ=0,\quad
 C_k:Z-\omega^kX=0\qquad(i,j,k\bmod3).
\]
The three same-group intersections are the coordinate vertices. The other
nine are
\[
 P_{ij}=[\omega^{i+j}:\omega^j:1],
\]
each incident to exactly
\(A_i,B_j,C_{-i-j}\). All twelve are ordinary triples, and they exhaust
the \(12\binom32=\binom92\) pairs. In particular there are no further
intersections. This is the dual Hesse arrangement of \cite{DvS}.

The infinity \(L=X+2Y+4Z=0\) misses the vertices and every \(P_{ij}\),
since \(|\omega^{i+j}+2\omega^j+4|\ge1\).
In the chart \(u=X/L,v=Y/L\), all nine affine directions are distinct.
Scaling \((u,v)=(x/t,y/t)\) gives the equation
\begin{equation}\label{F}
 F=(x^3-y^3)\bigl(64y^3-(t-x-2y)^3\bigr)
              \bigl((t-x-2y)^3-64x^3\bigr)=0.
\end{equation}
Its nine linear factors, indexed by \(q\in\{1,\omega,\omega^2\}\), are
\[
 x-qy,\qquad qx+(4+2q)y-qt,\qquad -(1+4q)x-2y+t.
\]
For \(t\ne0\) its singular set is exactly the scaled twelve-point set;
at zero the nine distinct directions meet only at the origin. Every fiber
is reduced. The coefficient of \(x^9\) is the constant \(-65\), so the
polynomial quotient is free of rank nine over \(\C[y,t]\) and flat over
\(\C[t]\). Directly in the analytic local ring at parameter \(t_0\),
put \(s=t-t_0\). If \(sg=Fh\), reduction modulo \(s\) in the domain
\(\C\{x-x_0,y-y_0\}\) shows \(h=sh_1\), since \(F_{t_0}\ne0\).
Cancellation gives \(g=Fh_1\). The quotient is torsion-free over the
base discrete valuation ring, hence flat. This persists on open restriction.

Each linear factor has nonzero \(x\)-coefficient, so its component plane
is a graph over \((y,t)\). The disjoint union of the nine planes maps
finitely to \(F=0\): on rings its algebra is the finite product of the
nine cyclic component quotients. The component ideals intersect in their
product, making this an injective extension with the same total fraction
ring. The product algebra is normal, finite and integral over the reduced
quotient, hence is its integral closure. At every fixed \(t\), including
zero, it restricts to the disjoint union of nine distinct smooth lines,
which is the normalization of that fiber. This is finite simultaneous
normalization, with normalized projection smooth. The total space \(F=0\)
itself is nonnormal. Analytic restriction preserves this description.

Use the open unit ball for the holomorphic family, and its closed ball for
boundary statements. The vertex chart coordinates are \((1,0),(0,1/2),(0,0)\).
At the other points, both coordinate numerators have modulus one and the
common denominator has modulus at least one. Hence every singular point
has norm at most \(\sqrt2|t|<1/2\). For any of the displayed lines
\(ax+by+ct=0\), its distance to the origin is
\(d(t)=|ct|/(|a|^2+|b|^2)^{1/2}\).
The three groups have bounds respectively \(0,|t|/\sqrt5,|t|/\sqrt{13}\),
all strictly below one. A closest point and a fixed unit direction are
\[
 p(t)=-\frac{ct(\bar a,\bar b)}{|a|^2+|b|^2},\qquad
 v=\frac{(-b,a)}{\sqrt{|a|^2+|b|^2}}.
\]
The sphere intersection is
\(p(t)+\sqrt{1-d(t)^2}\,e^{i\theta}v\), a transverse smoothly varying
circle. These nine circles are disjoint because all pair intersections
are strictly inside the half-radius ball. Thus they give a link isotopy
along any parameter path, and each normalized restricted line is a disk.
The central cone is transverse to every sphere about zero. Homogeneity
\(F(\rho x,\rho y,\rho t)=\rho^9F(x,y,t)\) gives the corresponding
radius-\(\rho\) representative over \(|t|<\rho/4\). No properness of
the open-ball projection is needed; the normalization map is finite.

For an ordinary \(r\)-fold point, transverse smooth branches give
\(\delta_r=\binom r2\), by the additivity formula
\(\delta(fg)=\delta(f)+\delta(g)+I(f,g)\).
For \(r\ge3\), one blowup has \(E^2=-1\), \(E\cdot C'=r\), and
\(K=E\): adjunction gives \((K+E)E=-2\).
Consequently
\[
 \rough_r=E\cdot(C'+2E)=r-2.
\]
A node already has normal crossings and rough value zero, in agreement
with this formula at \(r=2\). Applying it at the ordinary ninefold point
and all twelve triples proves both displayed equalities and the strict
inequality. In particular delta constancy is an equality of total local
invariants, with all singularities retained in the representative.
\end{proof}

\paragraph{Classical Fermat extension.}
For every integer \(n\ge3\), the classical arrangement
\[
 (X^n-Y^n)(Y^n-Z^n)(Z^n-X^n)=0
\]
has \(3n\) lines, three ordinary
\(n\)-fold vertices and \(n^2\) other ordinary triples
\cite[p.~119]{H}. Indeed the latter are
\([\zeta^{i+j}:\zeta^j:1]\); exactly one line of each group passes
through each, and same-group pairs meet at the vertices. These exhaust
all pairs. The same infinity misses every intersection. Scaling its
affine factors gives, by the proof above, finite simultaneous normalization
and a sufficiently small common transverse ball. Therefore
\[
 3\binom n2+3n^2=\binom{3n}2,\qquad
 \sum_p\rough(C_t,p)-\rough(C_0,0)
   =3(n-2)+n^2-(3n-2)=n^2-4.
\]
The excesses are unbounded. This conclusion follows from the all-\(n\)
incidence argument, rather than a finite computational range. At \(n=2\)
the vertices are nodes and the excess is zero.

\paragraph{Limits of the conclusion.}
The sums here run over all moving singularities; no counterexample to a
marked-section assertion is claimed. The central curve has nine branches.
Borodzik's published Definition~2.1 and Conjecture~3.8
\cite[pp.~575,577]{B} require a parametric deformation to have a
\emph{unibranched} central fiber as well as rational normalization.
The present family fails that hypothesis and does not resolve that
conjecture. We use precisely rough \(K(K+D)\), without substituting a
different multibranch convention or a fine invariant. For these ordinary
points \(K+D\) is relatively nef, so the negative part is zero; this does
not identify fine-invariant conventions on other singularities.

\paragraph{Verification and review status.}
The companion source and verification archive contains the unchanged
original 28-file packet, independent exact arithmetic controls, reviewed
provenance records, and a verifier using externally anchored hashes and
isolated nonoptimized replay. Computation checks incidence, factorization,
invariant arithmetic and boundary estimates; it supplements the analytic
proof and proves neither priority nor a universal assertion by sampling.
Raw third-party source bodies are excluded. The original recovery record
says ``1 verified/5''; the earlier attempt count is unavailable.
AI tools were used extensively in the research, writing and independent
adversarial checking. Those checks are not independent human peer review.
This preprint is unrefereed.

\clearpage
\begin{thebibliography}{9}
\raggedright
\bibitem{OWR} M.~Borodzik, On deformations of plane curve singularities,
in \emph{Singularities}, Oberwolfach Reports \textbf{6} (2009), no.~3,
2446--2448. Report published June~30, 2010.
\href{https://doi.org/10.4171/OWR/2009/43}{doi:10.4171/OWR/2009/43}.
\bibitem{B} M.~Borodzik, Deformations of singularities of plane curves:
topological approach, Osaka J. Math. \textbf{51} (2014), no.~3, 573--583.
\href{https://doi.org/10.18910/50809}{doi:10.18910/50809}.
\bibitem{DvS} T.~de Jong and D.~van Straten, Deformation theory of
sandwiched singularities, Duke Math. J. \textbf{95} (1998), no.~3,
451--522. \href{https://doi.org/10.1215/S0012-7094-98-09513-8}{doi:10.1215/S0012-7094-98-09513-8}.
\bibitem{PS} O.~Plamenevskaya and L.~Starkston, Unexpected Stein fillings,
rational surface singularities and plane curve arrangements,
Geom. Topol. \textbf{27} (2023), no.~3, 1083--1202.
\href{https://doi.org/10.2140/gt.2023.27.1083}{doi:10.2140/gt.2023.27.1083}.
\bibitem{H} F.~Hirzebruch, Arrangements of lines and algebraic surfaces,
in \emph{Arithmetic and Geometry}, Vol.~II, Progress in Mathematics 36,
Birkh\"auser, 1983, 113--140.
\href{https://doi.org/10.1007/978-1-4757-9286-7_7}{doi:10.1007/978-1-4757-9286-7\_7}.
\end{thebibliography}
\end{document}
